\documentclass[11pt,a4paper]{article}
\usepackage{../pelm}
\begin{document}
\answerpaper{Quiz 5}{Weeks 9--12: Agents, Design, Domain Adaptation, and Security}

\begin{enumerate}[leftmargin=*, itemsep=0.5em, label=\textbf{\arabic*.}]
\key{C}{With a chatbot you read every output and correct it unconsciously. An agent deletes that
  free quality control, and most people do not notice they were providing it.}
\key{C}{$0.95^{10} \approx 0.60$. Every component is excellent and the system fails about two times
  in five.}
\key{B}{Errors are inherited rather than averaged away. Step two performs its work on
  possibly-wrong input, so a mistake early is carried by everything downstream.}
\key{B}{Silent failure --- the characteristic failure of anything multi-step. Nothing errors, and the
  output is indistinguishable from a correct one.}
\key{C}{Narrow tools is a wall; everything else is a request. It cannot misuse a capability it was
  never given, regardless of what it decides.}
\key{B}{Before the irreversible action, not at the end. Identify the irreversible step first, then
  place the human relative to it.}
\key{B}{``A student misses a resit because they were told the wrong date'' immediately rules out
  designs that ``users could get wrong information'' permits.}
\key{C}{A Domain Pack makes output usable and more persuasive. It does nothing for correctness ---
  and by making wrong content look professional, it can increase over-trust.}
\key{B}{You did nothing wrong and never saw the malicious text. Direct injection is someone
  attacking their own tool; indirect injection is someone attacking you.}
\key{B}{Instructions and data share one undifferentiated channel. Nothing in the context carries a
  label marking it trustworthy, which is why this is hard to fix rather than merely neglected.}
\end{enumerate}
\end{document}
